\slajdid{10-s023}
\begin{frame}[label=10-s023]
\gusttekst\setlength{\abovedisplayskip}{3pt}\setlength{\belowdisplayskip}{3pt}\setlength{\abovedisplayshortskip}{2pt}\setlength{\belowdisplayshortskip}{2pt}Da bi odredili $I_{OH}$ smatraćemo da je na izlazu napon logičke jedinice. Tada je tranzistor zakočen pa je po izlaznoj konturi
\[I_O=-\frac{V_{CC}-V_O}{R_C}\]
Znak minus je posledica usvojenog referentog smera izlazne struje i on pokazuje ono što smo u elementima analize rekli da će logička kola za stanje logičke jedinice na izlazu biti izvori struje, odnosno da će pravi smer struje biti „iz kola“. I kao što smo u elementima analize rekli mogli bi sada u jednačini da zamenimo $V_O=V_{IH}$, pošto će naredno kolo taj napon i dalje shvatati kao logičku jedinicu. Ali u tom slučaju nismo ostavili prostor da sme da se pojavi bilo kakav šum na tom naponu. Zato definišemo minimalan napon logičke jedinice na izlazu $V_{OHmin}$ tako da je $V_{OHmin}>V_{IH}$. U tom slučaju
\[I_{Omax}=I_{OH}=-\frac{V_{CC}-V_{OHmin}}{R_C}\]
\end{frame}
